% Lamina 8: eleccion del punto de operacion del detector final (Ens-Vb-X).
% 2 x 2: columnas = banco | despliegue; filas = recall (arriba) | precision (abajo).
% Banco, p99 y p99.9: results/linea_base/rev24/metricas_banco.csv y falsas_alarmas_banco.csv
%   (Ens-Vb-X0.2596, E1; dias al año = fa x 365). Alineado con la lamina 9 el 2026-11-23.
% Banco, p95 y p98: corrida del registro 24-p98 con p95 agregado (metricas_e1.csv y
%   falsas_alarmas.csv, 2026-11-24), misma mascara que rev24 (23 005 limpias, 1 777 atacadas).
%   Reemplaza a percentiles_banco.csv (mascara comun, no estaba en el repo).
% Despliegue: revision 26, datos validados de SIMAT, unidad A (las mismas estacion-dias
%   del examen original): despliegue_validado_metricas.csv y _falsas_alarmas.csv.
% Umbrales de 2023 sin ataques para p95, p98, p99 y p99.9.
% Eje x: 1 = p95, 2 = p98, 3 = p99, 4 = p99.9 (de menos a mas estricto).
% Eje x: umbral (percentil de 2023 sin ataques). Debajo, en un recuadro aparte (complemento de la
% grafica, no parte de ella): dias con falsa alarma al año por estacion de cada percentil (2026-11-23).
% Tono = tamaño del ataque (claro = chico). Linea solida con circulo = anomalia proporcional
% (delta %); linea punteada con cuadro = error fijo (ppb). Recall arriba, precision abajo.
% Franja crema = p99, el elegido. (Estilos cambiados el 2026-11-23 a pedido del tesista.)
\begin{tikzpicture}[
    recuadroFA/.style={draw=tenue!60, line width=0.4pt, rounded corners=2pt},
    numFA/.style={font=\sffamily\scriptsize, text=tenue},
    textoFA/.style={font=\sffamily\tiny, text=tenue},
  ]
\begin{groupplot}[
    group style={group size=2 by 2, horizontal sep=0.45cm, vertical sep=0.3cm,
                 xticklabels at=edge bottom, yticklabels at=edge left},
    ejes tesis,
    scale only axis, width=5.4cm, height=1.6cm,
    xmin=0.6, xmax=4.4, xtick={1,2,3,4},
    xticklabel style={align=center},
    ymin=0, ymax=1.05, ytick={0,0.5,1},
    every axis plot/.append style={mark size=1.6pt},
    title style={font=\sffamily\small\bfseries, text=noche, yshift=-3pt},
    xlabel style={font=\sffamily\scriptsize, text=texto, yshift=3pt},
    clip=false,
  ]
  % ---------------- Banco · recall ----------------
  \nextgroupplot[title={Banco}, ylabel={Recall},
                 legend to name=leyendaPO, legend columns=4,
                 legend style={/tikz/every even column/.append style={column sep=0.35cm}}]
  \addplot[forget plot, fill=resalte, draw=none, mark=none] coordinates {(2.65,0) (2.65,1.05) (3.35,1.05) (3.35,0)} \closedcycle;  % franja p99
  \addlegendimage{empty legend}\addlegendentry{\textbf{Proporcional} (sólida):}
  \addplot[mag chica, fam pct, met recall] coordinates {(1,0.45) (2,0.29) (3,0.20) (4,0.03)}; \addlegendentry{$\delta$\,20\,\%}
  \addplot[mag media, fam pct, met recall] coordinates {(1,0.92) (2,0.85) (3,0.80) (4,0.40)}; \addlegendentry{$\delta$\,50\,\%}
  \addplot[mag grande, fam pct, met recall] coordinates {(1,0.98) (2,0.95) (3,0.93) (4,0.64)}; \addlegendentry{$\delta$\,100\,\%}
  \addlegendimage{empty legend}\addlegendentry{\textbf{Fijo} (punteada):}
  \addplot[mag chica, fam ppb, met recall] coordinates {(1,0.49) (2,0.29) (3,0.23) (4,0.05)}; \addlegendentry{+10 ppb}
  \addplot[mag media, fam ppb, met recall] coordinates {(1,0.95) (2,0.88) (3,0.83) (4,0.44)}; \addlegendentry{+20 ppb}
  \addplot[mag grande, fam ppb, met recall] coordinates {(1,0.99) (2,0.97) (3,0.96) (4,0.78)}; \addlegendentry{+40 ppb}
  % ---------------- Despliegue · recall ----------------
  \nextgroupplot[title={Despliegue}]
  \addplot[forget plot, fill=resalte, draw=none, mark=none] coordinates {(2.65,0) (2.65,1.05) (3.35,1.05) (3.35,0)} \closedcycle;  % franja p99
  \addplot[mag chica, fam pct, met recall] coordinates {(1,0.45) (2,0.30) (3,0.25) (4,0.03)};  % $\delta$\,20\,\%
  \addplot[mag media, fam pct, met recall] coordinates {(1,0.99) (2,0.97) (3,0.94) (4,0.53)};  % $\delta$\,50\,\%
  \addplot[mag grande, fam pct, met recall] coordinates {(1,1.00) (2,0.99) (3,0.99) (4,0.73)};  % $\delta$\,100\,\%
  \addplot[mag chica, fam ppb, met recall] coordinates {(1,0.55) (2,0.39) (3,0.33) (4,0.06)};  % 10 ppb
  \addplot[mag media, fam ppb, met recall] coordinates {(1,0.97) (2,0.96) (3,0.90) (4,0.53)};  % 20 ppb
  \addplot[mag grande, fam ppb, met recall] coordinates {(1,1.00) (2,1.00) (3,1.00) (4,0.86)};  % 40 ppb
  % ---------------- Banco · precision ----------------
  \nextgroupplot[ylabel={Precisión}, xlabel={Umbral: percentil de 2023 sin ataques},
                 xticklabels={p95, p98, p99, p99.9}]
  % recuadro aparte (complemento, no parte de la grafica): dias con falsa alarma por percentil
  \draw[recuadroFA] ([yshift=-0.93cm]rel axis cs:0,0) rectangle ([yshift=-1.55cm]rel axis cs:1,0);
  \node[numFA] at ([yshift=-1.09cm]axis cs:1,0) {18.5};
  \node[numFA] at ([yshift=-1.09cm]axis cs:2,0) {8.4};
  \node[numFA] at ([yshift=-1.09cm]axis cs:3,0) {4.7};
  \node[numFA] at ([yshift=-1.09cm]axis cs:4,0) {0.7};
  \node[textoFA] at ([yshift=-1.39cm]rel axis cs:0.5,0) {Días con falsa alarma al año por estación};
  \addplot[forget plot, fill=resalte, draw=none, mark=none] coordinates {(2.65,0) (2.65,1.05) (3.35,1.05) (3.35,0)} \closedcycle;  % franja p99
  \node[font=\sffamily\scriptsize\bfseries, text=acento, anchor=north] at (axis cs:3,1.03) {elegido};
  \addplot[mag chica, fam pct, met prec] coordinates {(1,0.06) (2,0.09) (3,0.11) (4,0.13)};  % $\delta$\,20\,\%
  \addplot[mag media, fam pct, met prec] coordinates {(1,0.12) (2,0.22) (3,0.31) (4,0.64)};  % $\delta$\,50\,\%
  \addplot[mag grande, fam pct, met prec] coordinates {(1,0.12) (2,0.22) (3,0.32) (4,0.73)};  % $\delta$\,100\,\%
  \addplot[mag chica, fam ppb, met prec] coordinates {(1,0.07) (2,0.09) (3,0.12) (4,0.17)};  % 10 ppb
  \addplot[mag media, fam ppb, met prec] coordinates {(1,0.12) (2,0.22) (3,0.32) (4,0.65)};  % 20 ppb
  \addplot[mag grande, fam ppb, met prec] coordinates {(1,0.12) (2,0.23) (3,0.33) (4,0.77)};  % 40 ppb
  % ---------------- Despliegue · precision ----------------
  \nextgroupplot[xlabel={Umbral: percentil de 2023 sin ataques},
                 xticklabels={p95, p98, p99, p99.9}]
  % recuadro aparte (complemento, no parte de la grafica): dias con falsa alarma por percentil
  \draw[recuadroFA] ([yshift=-0.93cm]rel axis cs:0,0) rectangle ([yshift=-1.55cm]rel axis cs:1,0);
  \node[numFA] at ([yshift=-1.09cm]axis cs:1,0) {5.4};
  \node[numFA] at ([yshift=-1.09cm]axis cs:2,0) {1.3};
  \node[numFA] at ([yshift=-1.09cm]axis cs:3,0) {1.1};
  \node[numFA] at ([yshift=-1.09cm]axis cs:4,0) {0.0};
  \node[textoFA] at ([yshift=-1.39cm]rel axis cs:0.5,0) {Días con falsa alarma al año por estación};
  \addplot[forget plot, fill=resalte, draw=none, mark=none] coordinates {(2.65,0) (2.65,1.05) (3.35,1.05) (3.35,0)} \closedcycle;  % franja p99
  \addplot[mag chica, fam pct, met prec] coordinates {(1,0.25) (2,0.44) (3,0.46) (4,1.00)};  % $\delta$\,20\,\%
  \addplot[mag media, fam pct, met prec] coordinates {(1,0.39) (2,0.63) (3,0.67) (4,0.96)};  % $\delta$\,50\,\%
  \addplot[mag grande, fam pct, met prec] coordinates {(1,0.35) (2,0.59) (3,0.65) (4,0.96)};  % $\delta$\,100\,\%
  \addplot[mag chica, fam ppb, met prec] coordinates {(1,0.29) (2,0.50) (3,0.52) (4,1.00)};  % 10 ppb
  \addplot[mag media, fam ppb, met prec] coordinates {(1,0.38) (2,0.63) (3,0.67) (4,0.97)};  % 20 ppb
  \addplot[mag grande, fam ppb, met prec] coordinates {(1,0.36) (2,0.59) (3,0.65) (4,0.98)};  % 40 ppb
\end{groupplot}
\end{tikzpicture}

\smallskip
\ref{leyendaPO}   % leyenda comun (se llena en la segunda compilacion)